\documentclass[11pt]{article}

% ── Encoding ──
\usepackage[utf8]{inputenc}
\usepackage[T1]{fontenc}

% ── Layout ──
\usepackage{geometry}
\geometry{margin=1in}

% ── Math ──
\usepackage{amsmath}
\usepackage{amssymb}
\usepackage{amsthm}

% ── Figures and tables ──
\usepackage{graphicx}
\usepackage{booktabs}

% ── References ──
\usepackage[colorlinks=true,linkcolor=blue,citecolor=blue,urlcolor=blue]{hyperref}

% ── Theorem environments ──
\newtheorem{theorem}{Theorem}[section]
\newtheorem{lemma}[theorem]{Lemma}
\newtheorem{proposition}[theorem]{Proposition}
\newtheorem{corollary}[theorem]{Corollary}
\theoremstyle{definition}
\newtheorem{definition}[theorem]{Definition}
\theoremstyle{remark}
\newtheorem{remark}[theorem]{Remark}
\newtheorem{conjecture}[theorem]{Conjecture}
\newtheorem{question}[theorem]{Question}

\title{Universal Period-3 Orbits for Generalized H\'enon Maps over~$\mathbb{Q}$:\\
A Counterexample to Conjecture~4 of Berger et al.}

\author{Guilherme de Azevedo Silveira\\
\texttt{guilherme.silveira@gmail.com}}
\date{\today}

\begin{document}

\maketitle

\begin{abstract}
The generalized H\'enon map $f_{c,b}(x,y) = (y,\, y^2 + c + bx)$, with
$c, b \in \mathbb{Q}$, is a polynomial automorphism of the affine plane whose
arithmetic dynamics has been the subject of recent conjectures.
Conjecture~4 of Berger et al., originating from the
2023 BIRS workshop, asserts that for all but finitely many $\delta \in \mathbb{Q}$,
every $\mathbb{Q}$-rational periodic point of $f_{c,\delta}$ has period dividing~$2$.
We disprove this conjecture in the strongest possible way by exhibiting an explicit
algebraic family: for every $b \in \mathbb{Q} \setminus \{-1\}$, the map $f_{c,b}$
admits a primitive period-$3$ orbit over~$\mathbb{Q}$, with
$c = -(29b^2 + 38b + 29)/16$.
The proof is a direct algebraic verification of three polynomial identities,
formally verified in Lean~4 with Mathlib.
We also verify the sharpness of Conjecture~3 by exhibiting rational
orbits for all six allowed periods $\{1,2,3,4,6,8\}$, discover a period-$7$ orbit
showing that the period set depends on the sign of the Jacobian determinant,
and use genus analysis to reconcile the universal family with Faltings finiteness
for each fixed~$b$.
\end{abstract}

\medskip
\noindent\textbf{Keywords:} H\'enon map; rational periodic points; arithmetic dynamics;
counterexample; universal parametric family; dynatomic polynomial; Faltings finiteness;
BIRS 2023; formal verification; Lean~4

\medskip
\noindent\textbf{MSC 2020:} 37P05, 37P15, 14G05, 11G30


% ============================================================================
\section{Introduction}
\label{sec:intro}
% ============================================================================

Let $b \in \mathbb{Q}^*$ and $c \in \mathbb{Q}$. The \emph{generalized H\'enon map}
is the polynomial automorphism of $\mathbb{A}^2$:
%
\begin{equation}\label{eq:henon}
  f_{c,b}\colon (x,y) \longmapsto (y,\; y^2 + c + bx).
\end{equation}
%
The Jacobian determinant is $\det Df_{c,b} = -b$, so $f$ is area-preserving when
$|b| = 1$, orientation-reversing when $b > 0$, and orientation-preserving when $b < 0$.
% SOURCE: Df = [[0,1],[b,2y]], det = 0·2y - 1·b = -b

The arithmetic dynamics of~\eqref{eq:henon} --- the study of $\mathbb{Q}$-rational
periodic orbits (see Silverman~\cite{silverman2007} for general background and
Benedetto et al.~\cite{benedetto2019} for a recent survey) ---
was examined in the survey of open problems by Berger, Bedford,
and fourteen co-authors~\cite{berger2024}, arising from the 2023 BIRS workshop.
Two conjectures are central.

\begin{conjecture}[{Conjecture~3 of \cite{berger2024}}]\label{conj:3}
  Over $\mathbb{Q}$, the map $f_{c,b}$ with $b = +1$
  (i.e., $f_{c,\delta}$ with $\delta = -1$ in Berger et al.'s convention),
  namely $(x,y) \mapsto (y,\, y^2 + c + x)$,
  has no rational point of primitive period $N$ other than $N \in \{1, 2, 3, 4, 6, 8\}$.
\end{conjecture}

\begin{conjecture}[{Conjecture~4 of \cite{berger2024}}]\label{conj:4}
  For all but finitely many $\delta \in \mathbb{Q}$, the $\mathbb{Q}$-rational periodic
  points of $f_{c,\delta}$ (for any $c \in \mathbb{Q}$) have period dividing~$2$.
\end{conjecture}

Conjecture~\ref{conj:4} predicts a sharp dichotomy: generically, only fixed points
and period-$2$ cycles exist over~$\mathbb{Q}$, with at most finitely many exceptional
parameters admitting longer periods.  This prediction is natural by analogy with
the one-dimensional case, where Poonen conjectured that the quadratic map
$z \mapsto z^2 + c$ has no rational periodic point of period $\geq 4$
(see Flynn--Poonen--Schaefer~\cite{flynnpoonenschaefer1997}).
In the two-dimensional setting, however, the extra parameter~$b$ opens a degree of
freedom absent in the one-variable case.

Our main result exploits this extra parameter to disprove Conjecture~\ref{conj:4}
in the strongest possible way.

\begin{theorem}[Universal period-3 family]\label{thm:main}
  For every $b \in \mathbb{Q} \setminus \{-1\}$, set
  %
  \begin{equation}\label{eq:family}
    c = -\frac{29b^2 + 38b + 29}{16}, \quad
    \begin{cases}
      x_0 = -\dfrac{5b+7}{4},\\[6pt]
      x_1 = \dfrac{7b+5}{4},\\[6pt]
      x_2 = \dfrac{b-1}{4}.
    \end{cases}
  \end{equation}
  %
  Then $(x_0, x_1, x_2)$ is a primitive period-$3$ orbit of
  $f_{c,b}$ over~$\mathbb{Q}$.
\end{theorem}

\begin{corollary}\label{cor:conj4false}
  Every $b \in \mathbb{Q} \setminus \{-1\}$ admits a $\mathbb{Q}$-rational orbit
  of primitive period~$3$.  The ``exceptional set'' in Conjecture~\ref{conj:4} is not
  finite --- it is all of $\mathbb{Q} \setminus \{-1\}$.
  Conjecture~\ref{conj:4} is false.
\end{corollary}

\begin{remark}[Formal verification]\label{rem:lean}
  Theorem~\ref{thm:main}, Corollary~\ref{cor:conj4false}, and
  Proposition~\ref{prop:period7} have been formally verified in the
  Lean~4 proof assistant~\cite{lean4} using Mathlib, with zero \texttt{sorry}
  and zero custom axioms.  The formalization is available in the
  \texttt{lean/} directory of the accompanying code repository.
\end{remark}

At $b = -1$ (the orientation-preserving area-preserving case),
the family~\eqref{eq:family} degenerates to a fixed point (see
Remark~\ref{rem:degenerate}), but period-$3$ orbits exist via a different mechanism.
Thus period~$3$ is universal across \emph{all} of~$\mathbb{Q}$.

\smallskip
\noindent\textbf{Secondary results.}
\begin{enumerate}
  \item \textbf{Sharpness of Conjecture~\ref{conj:3}}
  (Section~\ref{sec:conj3}):
  We exhibit rational orbits for every allowed period $N \in \{1,2,3,4,6,8\}$ and
  explain why the denominator-pattern family $c = -(2q+1)/q^2$ produces periods
  $2q$ only for $q \leq 4$.

  \item \textbf{Period~$7$ and orientation dependence}
  (Section~\ref{sec:period7}):
  At $b = -1$, $c = -9/16$, we discover a period-$7$ orbit.
  The same~$c$ gives period~$8$ at $b = +1$, revealing that the
  period set depends on the sign of $\det Df = -b$.

  \item \textbf{Genus analysis and Faltings finiteness}
  (Section~\ref{sec:genus}):
  For fixed~$b$, the primitive period-$3$ dynatomic curve has genus~$\geq 9$,
  so Faltings' theorem gives finitely many $c \in \mathbb{Q}$ with period-$3$ orbits.
  The universal family avoids this constraint by varying~$b$, tracing a rational
  (genus-$0$) curve in the $(b,c)$-plane.
\end{enumerate}

\smallskip
\noindent\textbf{Organization.}
Section~\ref{sec:prelim} fixes notation.
Section~\ref{sec:main} proves Theorem~\ref{thm:main}.
Section~\ref{sec:evidence} presents the computational discovery process.
Section~\ref{sec:conj3} treats Conjecture~\ref{conj:3} sharpness and the
denominator-pattern family.
Section~\ref{sec:period7} discusses period~$7$ and orientation dependence.
Section~\ref{sec:genus} gives the genus analysis.
Section~\ref{sec:discussion} collects open questions.


% ============================================================================
\section{Preliminaries}
\label{sec:prelim}
% ============================================================================

\begin{definition}\label{def:periodicorbit}
  A \emph{periodic orbit of period~$n$} for~\eqref{eq:henon} is a sequence
  $(x_0, x_1, \ldots, x_{n-1}) \in \mathbb{Q}^n$ satisfying the recurrence
  %
  \begin{equation}\label{eq:recurrence}
    x_{i+1} = x_i^2 + c + b\, x_{i-1} \qquad (i = 1, \ldots, n-1),
  \end{equation}
  %
  with closure conditions $x_n = x_0$ and $x_{n+1} = x_1$.
  The period is \emph{primitive} if no proper divisor of~$n$ is also a period.
\end{definition}

Equivalently, $(x_0, x_1)$ is a periodic point of $f_{c,b}$ in the usual
$(x,y)$-plane: the orbit under~$f$ is
$(x_0, x_1) \to (x_1, x_2) \to \cdots \to (x_{n-1}, x_0) \to (x_0, x_1)$.

\begin{remark}[Sign convention]\label{rem:convention}
  Berger et al.~\cite{berger2024} use the convention
  $f_{c,\delta}(x,y) = (y,\, y^2 + c - \delta x)$ with parameter~$\delta$.
  Our parameter~$b$ relates to their~$\delta$ via $b = -\delta$.
  In particular, Conjecture~\ref{conj:3} concerns $\delta = -1$, which
  corresponds to $b = +1$ in our convention (the orientation-reversing
  area-preserving case, $\det Df = -b = -1$).
  We state the conjectures using Berger et al.'s original notation
  and convert to our convention~$b$ for all proofs and computations.
\end{remark}

For period~$3$, the recurrence~\eqref{eq:recurrence} and closure conditions give
three equations in $(x_0, x_1, x_2, c)$ parametrized by~$b$:
%
\begin{align}
  x_2 &= x_1^2 + c + b\, x_0, \label{eq:p3a}\\
  x_0 &= x_2^2 + c + b\, x_1, \label{eq:p3b}\\
  x_1 &= x_0^2 + c + b\, x_2. \label{eq:p3c}
\end{align}


% ============================================================================
\section{Proof of Theorem~\ref{thm:main}}
\label{sec:main}
% ============================================================================

\begin{proof}
  Let $b \in \mathbb{Q}$, $b \neq -1$, and define $c, x_0, x_1, x_2$
  by~\eqref{eq:family}.  We verify
  equations~\eqref{eq:p3a}--\eqref{eq:p3c} as polynomial identities in~$b$.
  Each identity was independently verified using the computer algebra system
  SymPy~\cite{sympy}.
  % SOURCE: proof_full_chain.py — all three identities = 0 as polynomials in b

  \smallskip\noindent\textbf{Equation~\eqref{eq:p3a}:}
  $x_1^2 + c + b\,x_0 = x_2$.
  \begin{align*}
    x_1^2 + c + b\,x_0
    &= \frac{(7b+5)^2}{16} - \frac{29b^2+38b+29}{16} + b\!\cdot\!\frac{-(5b+7)}{4} \\[4pt]
    &= \frac{49b^2+70b+25 - 29b^2-38b-29}{16} - \frac{5b^2+7b}{4} \\[4pt]
    &= \frac{20b^2+32b-4}{16} - \frac{20b^2+28b}{16} \\[4pt]
    &= \frac{4b-4}{16} = \frac{b-1}{4} = x_2. \quad \checkmark
  \end{align*}
  % SOURCE: proof_full_chain.py Step 2: f(P₀)_y - x₂ = 0

  \smallskip\noindent\textbf{Equation~\eqref{eq:p3b}:}
  $x_2^2 + c + b\,x_1 = x_0$.
  \begin{align*}
    x_2^2 + c + b\,x_1
    &= \frac{(b-1)^2}{16} - \frac{29b^2+38b+29}{16} + b\!\cdot\!\frac{7b+5}{4} \\[4pt]
    &= \frac{b^2-2b+1-29b^2-38b-29}{16} + \frac{7b^2+5b}{4} \\[4pt]
    &= \frac{-28b^2-40b-28}{16} + \frac{28b^2+20b}{16} \\[4pt]
    &= \frac{-20b-28}{16} = \frac{-(5b+7)}{4} = x_0. \quad \checkmark
  \end{align*}
  % SOURCE: proof_full_chain.py Step 2: f²(P₀)_y - x₀ = 0

  \smallskip\noindent\textbf{Equation~\eqref{eq:p3c}:}
  $x_0^2 + c + b\,x_2 = x_1$.
  \begin{align*}
    x_0^2 + c + b\,x_2
    &= \frac{(5b+7)^2}{16} - \frac{29b^2+38b+29}{16} + b\!\cdot\!\frac{b-1}{4} \\[4pt]
    &= \frac{25b^2+70b+49 - 29b^2-38b-29}{16} + \frac{b^2-b}{4} \\[4pt]
    &= \frac{-4b^2+32b+20}{16} + \frac{4b^2-4b}{16} \\[4pt]
    &= \frac{28b+20}{16} = \frac{7b+5}{4} = x_1. \quad \checkmark
  \end{align*}
  % SOURCE: proof_full_chain.py Step 2: f³(P₀)_y - x₁ = 0

  \smallskip\noindent\textbf{Primitivity.}
  The orbit is primitive (not a fixed point) because
  $x_1 - x_0 = \frac{7b+5}{4} + \frac{5b+7}{4} = 3(b+1)$,
  which is nonzero for $b \neq -1$.
  Since $3$ is prime, any periodic orbit with period dividing~$3$ and period $\neq 1$
  must have period exactly~$3$.
  % SOURCE: proof_full_chain.py Step 4: x₁ - x₀ = 3*(b + 1), zero only at b = -1
\end{proof}

\begin{remark}[All three orbit points are distinct]\label{rem:distinct}
  The three orbit points in $\mathbb{Q}^2$ are
  $P_0 = (x_0, x_1)$, $P_1 = (x_1, x_2)$, $P_2 = (x_2, x_0)$.
  One verifies that $P_i = P_j$ requires $b = -1$ in every case.
  Thus for $b \neq -1$, the three orbit points are pairwise distinct.
  % SOURCE: proof_full_chain.py Step 6: all degenerate sets = {-1}
\end{remark}

\begin{remark}[Stability: the universal orbit is always a saddle]\label{rem:saddle}
  The Jacobian of $f^3$ along the orbit has $\det(Df^3) = (-b)^3 = -b^3$ and
  trace
  % SOURCE: src/proof_stability.py Steps 3-4: det and trace verified symbolically
  \[
    \operatorname{tr}(Df^3) = -\frac{(b-1)(35b^2 + 62b + 35)}{8}.
  \]
  The characteristic polynomial is $\lambda^2 - \operatorname{tr}(Df^3)\lambda - b^3 = 0$,
  with discriminant $\Delta = \operatorname{tr}^2 + 4b^3$.
  One verifies that $\Delta = (b+1)^2 Q(b)$ where
  $Q(b) = (1225b^4 - 560b^3 - 1266b^2 - 560b + 1225)/64$ has no real roots and
  satisfies $Q(b) \geq 0.985$ for all $b \in \mathbb{R}$.
  % SOURCE: src/proof_stability.py Step 6: Q has no real roots, min ≈ 0.985
  Thus $\Delta \geq 0$ for all~$b$, with equality only at $b = -1$
  (where the orbit degenerates to a fixed point).
  The eigenvalues are always real: the orbit is a \emph{saddle} for every
  $b \in \mathbb{Q} \setminus \{-1\}$, never elliptic or attracting.

  At $b = +1$ (Conjecture~\ref{conj:3} case):
  $\operatorname{tr} = 0$ and $\det = -1$, giving eigenvalues $\lambda = \pm 1$ ---
  the orbit lies exactly on the stability boundary.
  At $b = 0$ (the quadratic family): $\lambda_1 = 35/8$, $\lambda_2 = 0$
  (degenerate, as $\det Df = 0$).
  % SOURCE: src/proof_stability.py Step 7: special cases b=0,±1 verified
  % SOURCE: src/proof_stability.py Step 8: 20 numerical spot-checks all pass
\end{remark}

\begin{remark}[Degeneration at $b = -1$]\label{rem:degenerate}
  At $b = -1$ (the orientation-preserving area-preserving case,
  $\det Df = +1$),
  the family~\eqref{eq:family} gives $x_0 = x_1 = x_2 = -1/2$ and $c = -5/4$,
  which is a fixed point of~$f_{c,b}$.
  However, $b = -1$ does admit period-$3$ orbits through a different algebraic
  mechanism: the dynatomic polynomial $\Phi_3(x_0,c;\, b\!=\!-1)$ factors as
  $(c + x_0^2 + 1)^2(c + x_0^2 + 2x_0 + 2)$,
  and both linear-in-$c$ factors yield parametric families.
  For example, the factor $c = -(x_0^2 + 1)$ gives the primitive
  period-$3$ orbit $(0,\, -1,\, 0)$ at $c = -1$,
  the orbit $(1,\, -2,\, 1)$ at $c = -2$, and infinitely many others.
  % SOURCE: Φ₃ factorization at b=-1, verified computationally
  % (Corrected from earlier version that cited exp001e c=-7/9: that is a
  %  period-6 orbit at b=+1, not period-3 at b=-1. Corrected by ACT-01 audit.)
\end{remark}

\begin{remark}[Odd integer specialization]\label{rem:odd}
  When $b = 2k+1$ for $k \in \mathbb{Z}$ (i.e., $b$ is any odd integer),
  the orbit coordinates simplify to
  $x_0 = -(5k+6)/2$, $x_1 = (7k+6)/2$, $x_2 = k/2$,
  $c = -(29k^2+48k+24)/4$,
  giving orbits with especially small height (integer or half-integer coordinates).
  This special case was the intermediate discovery that led to the
  universal family~\eqref{eq:family}.
  % SOURCE: prove_family.py, exp002i CONCLUSIONS
  % (Corrected: formulas are for b=+(2k+1), not b=-(2k+1). The original
  %  exp002i used δ-convention where δ=2k+1; in b-convention, b=-δ=-(2k+1).
  %  But the orbit formulas were derived for the δ value, which equals b=+(2k+1)
  %  in the paper's convention. Corrected by ACT-01 round-3 audit.)
\end{remark}


% ============================================================================
\section{Discovery Process and Computational Evidence}
\label{sec:evidence}
% ============================================================================

The universal family~\eqref{eq:family} was found through a sequence of
computational explorations using exact rational arithmetic
(Python \texttt{Fraction} class, arbitrary precision).
We describe the discovery path, both to explain how a result this elementary
was not previously known and to illustrate the power of computational
exploration in arithmetic dynamics.

\subsection{Initial search: apparent rarity of exceptions}\label{sec:initial}

We tested $895$ values of $\delta = p/q$ with $q \leq 24$ and $|p| \leq 30$,
searching for period-$3$ orbits among $(c, x_0, x_1)$ with small denominators
($\leq 2$) and $|c| \leq 10$, $|x_0|, |x_1| \leq 5$.
% SOURCE: exp002 CONCLUSIONS: 895 delta tested

Only $\delta = \pm 1$ appeared exceptional.  This was consistent with
Conjecture~\ref{conj:4} and suggested that the exceptional set might indeed
be $\{-1, +1\}$ (the area-preserving cases).

\subsection{Algebraic expansion: odd integers}

Algebraic analysis of the period-$3$ conditions --- solving the discriminant
of the quadratic in~$x_2$ from equation~\eqref{eq:p3a} --- revealed that
all odd integer values of~$\delta$ admit period-$3$ orbits, but with
$c$~growing as $O(\delta^2)$, exceeding the initial search range.
% SOURCE: exp002g CONCLUSIONS: odd integers found

The parametric family for odd~$\delta$ (Remark~\ref{rem:odd}) was then
proved symbolically.

\subsection{Even integers and beyond}

Extending the search range to $|c| \leq 500$ with orbit denominators up to~$4$
revealed that \emph{all} nonzero integers~$\delta$ admit period-$3$ orbits,
not just odd ones.  Even integers use orbits with denominator~$4$
(quarter-integer coordinates), which were invisible to the earlier search.
% SOURCE: exp002h CONCLUSIONS: ALL integers have period-3

\subsection{The universal family}

Reparametrizing the odd-integer family as a continuous function of~$b$ and
verifying symbolically, we discovered that the three period-$3$
equations~\eqref{eq:p3a}--\eqref{eq:p3c} vanish \emph{identically as
polynomials in~$b$}, not just at integer points.
The family works for every $b \in \mathbb{Q} \setminus \{-1\}$.
% SOURCE: unified_family_d.py — all three = 0 as polynomials in b

In particular, the apparent ``false negatives'' at half-integer values of~$b$
were caused by insufficient search range: the orbit coordinates have
denominator~$8$ when $b$ is a half-integer, well beyond the search bound
of denominator~$4$.
% SOURCE: exp002j CONCLUSIONS: half-integer orbits with x_denom=8

\subsection{Table of examples}

Table~\ref{tab:orbits} shows the family~\eqref{eq:family} for selected values
of~$b$, illustrating the range of orbit heights.

\begin{table}[ht]
  \centering
  \caption{Primitive period-$3$ orbits from the universal family~\eqref{eq:family}.
    All orbits verified by exact rational arithmetic.}
  \label{tab:orbits}
  \begin{tabular}{@{} rrrrrl @{}}
    \toprule
    $b$ & $c$ & $x_0$ & $x_1$ & $x_2$ & Note \\
    \midrule
    $-10$ & $-2549/16$ & $43/4$   & $-65/4$ & $-11/4$ & $b \in \mathbb{Z}$ \\
    $-5$  & $-141/4$   & $9/2$    & $-15/2$ & $-3/2$  & $b \in \mathbb{Z}$ \\
    $-2$  & $-69/16$   & $3/4$    & $-9/4$  & $-3/4$  & $b \in \mathbb{Z}$ \\
    $0$   & $-29/16$   & $-7/4$   & $5/4$   & $-1/4$  & Quadratic map \\
    $1$   & $-6$       & $-3$     & $3$     & $0$     & Conj.~3 case \\
    $2$   & $-221/16$  & $-17/4$  & $19/4$  & $1/4$   & $b \in \mathbb{Z}$ \\
    $5$   & $-59$      & $-8$     & $10$    & $1$     & $b \in \mathbb{Z}$ \\
    $10$  & $-3309/16$ & $-57/4$  & $75/4$  & $9/4$   & $b \in \mathbb{Z}$ \\
    \midrule
    $1/2$ & $-221/64$  & $-19/8$  & $17/8$  & $-1/8$  & Denom.~$8$ \\
    $3/7$ & $-155/49$  & $-16/7$  & $2$     & $-1/7$  & Denom.~$7$ \\
    $-17/5$ & $-1469/100$ & $5/2$ & $-47/10$ & $-11/10$ & Large height \\
    \bottomrule
  \end{tabular}
\end{table}
% SOURCE: proof_full_chain.py Step 7 (spot-check) and unified_family_d.py verification


\subsection{Why was this not found before?}\label{sec:whynew}

Several factors explain why the universal family~\eqref{eq:family}, despite being
an elementary algebraic identity, was not previously identified:
\begin{enumerate}
  \item \textbf{Recency.}  Conjecture~\ref{conj:4} was formulated at the 2023
  BIRS workshop and published in 2024~\cite{berger2024}.  The research community
  has had limited time to investigate it.

  \item \textbf{Theoretical misdirection.}  The standard tools in arithmetic
  dynamics --- dynatomic polynomials, canonical heights, Faltings' theorem ---
  prove that for each \emph{fixed}~$b$, only finitely many $c \in \mathbb{Q}$
  give period-$3$ orbits (Section~\ref{sec:genus}).  This finiteness, while
  correct, concerns the wrong fiber: it does not preclude an algebraic family
  parametrized by~$b$ itself.

  \item \textbf{One-dimensional analogy.}  In the one-variable case
  $z \mapsto z^2 + c$, rational periodic orbits are extremely constrained
  (Poonen's conjecture, \cite{flynnpoonenschaefer1997}).
  The extra parameter~$b$ in the H\'enon setting provides a degree of freedom
  absent in one dimension.

  \item \textbf{Computational gap.}  Discovering the family required
  systematic search with exact arithmetic followed by algebraic
  reparametrization --- a workflow more natural in experimental mathematics
  than in the theoretical arithmetic dynamics community.
\end{enumerate}


% ============================================================================
\section{Sharpness of Conjecture~\ref{conj:3} and the Denominator-Pattern Family}
\label{sec:conj3}
% ============================================================================

Conjecture~\ref{conj:3} asserts that the map $f_{c,b}$ with $b = +1$,
namely $(x,y) \mapsto (y, y^2+c+x)$,
admits $\mathbb{Q}$-rational orbits only for periods in $\{1,2,3,4,6,8\}$.
We verify that this conjecture, if true, is sharp.

\subsection{Orbits for each allowed period}

Table~\ref{tab:conj3} exhibits explicit rational orbits for all six allowed periods.

\begin{table}[ht]
  \centering
  \caption{Rational orbits of $f_{c,b}$ with $b = +1$ (i.e., $(x,y) \mapsto (y,\, y^2+c+x)$)
    for every period in $\{1,2,3,4,6,8\}$.}
  \label{tab:conj3}
  \begin{tabular}{@{} rlll @{}}
    \toprule
    Period $N$ & $c$ & Orbit $(x_0, x_1, \ldots)$ & Method \\
    \midrule
    $1$ & $-1/4$ & $(-1/2)$ & $c = -x_0^2$ \\
    $2$ & $-4$   & $(-2,\, 2)$ & $c = -x_0^2$, $x_1 = -x_0$ \\
    $3$ & $-6$   & $(-3,\, 3,\, 0)$ & $\Phi_3$ cubic \\
    $4$ & $-5/4$ & $(-3/2,\, 3/2,\, -1/2,\, 1/2)$ & $c = -(x_0^2 + 2x_0 + 2)$ \\
    $6$ & $-7/9$ & $(-1/3,\, 2/3,\, -2/3,\, 1/3,\, -4/3,\, 4/3)$ & $\Phi_6$ cubic factor \\
    $8$ & $-9/16$ & $(-5/4,\, 5/4,\, -1/4,\, 3/4,$ & Denom.\ pattern \\
        &         & $\phantom{(}-1/4,\, 1/4,\, -3/4,\, 1/4)$ & \\
    \bottomrule
  \end{tabular}
\end{table}
% SOURCE: exp001 (N=1,2,3,4), exp001e (N=6, c=-7/9), exp001f (N=8, c=-9/16)

Periods~$6$ and~$8$ were the most difficult to find.  The period-$6$ orbit was
obtained by factoring the dynatomic polynomial $\Phi_6(x,c;-1)$ into five
irreducible components (four cubics and one factor of degree~$18$) and finding
rational points on one of the cubic factors.
% SOURCE: exp001e CONCLUSIONS: Phi₆ factored, 4 cubics + degree-18
The period-$8$ orbit came from the following construction.

\subsection{The denominator-pattern family}\label{sec:denomfamily}

Setting $c = -(2q+1)/q^2$ and $x_0 = -(q+1)/q$, $x_1 = (q+1)/q$ for
$q \in \{2, 3, 4\}$ produces orbits of period $2q$ whose coordinates all
have denominator exactly~$q$:

\begin{center}
\begin{tabular}{@{} rrl @{}}
  \toprule
  $q$ & Period & Status \\
  \midrule
  $2$ & $4$ & $\checkmark$ \\
  $3$ & $6$ & $\checkmark$ \\
  $4$ & $8$ & $\checkmark$ \\
  $5$ & --- & Fails \\
  \bottomrule
\end{tabular}
\end{center}
% SOURCE: exp001f CONCLUSIONS: family c=-(2q+1)/q² for q=2,3,4

\begin{proposition}[Denominator preservation]\label{prop:denom}
  Let $c = -(2q+1)/q^2$ with $x_0 = -(q+1)/q$ and $x_1 = (q+1)/q$.
  Write $x_n = a_n/q$.  Then all orbit elements have denominator exactly~$q$
  if and only if $a_n^2 \equiv 1 \pmod{q}$ for all~$n$.
  The numerator recurrence is
  $a_{n+1} = (a_n^2 - 1)/q + a_{n-1} - 2$.
\end{proposition}
% SOURCE: exp001g CONCLUSIONS: denominator preservation theorem

The family fails at $q = 5$ because the numerator sequence begins
$a_0 = -6,\, a_1 = 6,\, a_2 = -1,\, a_3 = 4,\, a_4 = 0$,
and $0^2 \not\equiv 1 \pmod{5}$, causing the denominator to cascade from~$5$
to~$25$, $5^4$, $5^8$, \ldots
% SOURCE: exp001g CONCLUSIONS: q=5 fails at a₄=0

More generally, the fraction of residues~$a$ with $a^2 \equiv 1 \pmod{q}$
is at most $2/q$ for $q$ prime, which tends to zero.  The orbit must remain
in this ``safe set'' for~$2q$ steps, making success increasingly improbable
as~$q$ grows.  This provides an arithmetic explanation for why the family
cannot extend to period~$10$ and beyond, consistent with Conjecture~\ref{conj:3}.
% SOURCE: exp001g CONCLUSIONS: residue statistics, 2/q fraction

We also searched $3.3 \times 10^7$ parameter triples for forbidden periods
$N \in \{5, 7, 9, 10, 11, 12\}$, finding none.
% SOURCE: exp001 CONCLUSIONS: ~40M exact tests, periods 5,7,9,10,11,12 all 0


% ============================================================================
\section{Period~7 and Orientation Dependence}
\label{sec:period7}
% ============================================================================

\begin{proposition}\label{prop:period7}
  At $b = -1$ and $c = -9/16$, the map $f_{c,b}$ with $b = -1$
  (orientation-preserving, $\det Df = +1$)
  admits a primitive period-$7$
  orbit over~$\mathbb{Q}$.  All seven orbit points are distinct and have
  coordinates with denominator~$4$.
\end{proposition}
% SOURCE: exp004 CONCLUSIONS: period-7 orbit verified at b=-1, c=-9/16, 7 distinct states

The same value $c = -9/16$ gives period~$8$ at $b = +1$
(Table~\ref{tab:conj3}).  This shared parameter reveals a striking
orientation dependence:

\begin{center}
\begin{tabular}{@{} rcl @{}}
  \toprule
  $b$ & $\det Df = -b$ & Known periods over~$\mathbb{Q}$ \\
  \midrule
  $+1$ & $-1$ (orient.-reversing) & $\{1, 2, 3, 4, 6, 8\}$ (Conj.~\ref{conj:3}) \\
  $-1$ & $+1$ (orient.-preserving) & $\{1, 3, 4, 7, \ldots\}$ \\
  \bottomrule
\end{tabular}
\end{center}
% SOURCE: exp004 CONCLUSIONS: period sets for b=±1

Period~$7$ exists for the orientation-preserving map ($b = -1$) but not
for the orientation-reversing one ($b = +1$, where Conjecture~\ref{conj:3}
forbids it).

\begin{remark}
  We searched for period~$5$ at $b \in \{-1, +1, \pm 2, \pm 3, \pm 1/2\}$
  with $c \in [-100, 100]$ and orbit denominators up to~$4$, finding no
  period-$5$ orbits at any~$b$.  This is consistent with Conjecture~\ref{conj:3}
  for $b = -1$ and suggests that period~$5$ may be universally absent.
  % SOURCE: exp004 CONCLUSIONS: period 5 NONE found at any b
\end{remark}

\begin{remark}
  This result partially addresses Question~45 of~\cite{berger2024},
  which asks for infinite families of maps with rational cycles of length
  $\geq d + 3 = 5$.  We exhibit individual examples (period~$7$ at $b = -1$,
  period~$8$ at $b = +1$), but no infinite \emph{family} of length $\geq 5$
  was found.
  % SOURCE: exp004 CONCLUSIONS: Q45 partially answered
\end{remark}


% ============================================================================
\section{Genus Analysis and Faltings Finiteness}
\label{sec:genus}
% ============================================================================

Theorem~\ref{thm:main} shows that period-$3$ orbits exist for every
$b \in \mathbb{Q} \setminus \{-1\}$.  How is this compatible with the classical
finiteness results in arithmetic geometry?

\subsection{The dynatomic curve for fixed b}

Fix $b = +1$ (the Conjecture~\ref{conj:3} case, $\delta = -1$).  The period-$3$ conditions
\eqref{eq:p3a}--\eqref{eq:p3c} can be reduced (by eliminating $x_1$ and $x_2$)
to a polynomial equation $F(x_0, c) = 0$.  After factoring out the fixed-point
curve $c + x_0^2 = 0$, the primitive period-$3$ component is an irreducible
polynomial of degree~$6$ in~$(x_0, c)$:
% SOURCE: exp006 CONCLUSIONS: resultant degree 6, irreducible
%
\begin{equation}\label{eq:primitive3}
  x_0^6 + 3x_0^4 c + 3x_0^2 c^2 + 4x_0^2 c - 3x_0^2 + c^3 + 4c^2 - 11c + 6 = 0.
\end{equation}
% SOURCE: exp006 CONCLUSIONS: degree-6 factor explicitly

This plane curve has a singular point at $(0, 1)$ (an acnode), giving
%
\begin{equation}\label{eq:genus}
  g \geq \frac{(6-1)(6-2)}{2} - 1 = 9.
\end{equation}
% SOURCE: exp006 CONCLUSIONS: genus ≥ 9, singular at (0,1)

By Faltings' theorem~\cite{faltings1983}, a curve of genus $\geq 2$ over
$\mathbb{Q}$ has finitely many rational points.  Therefore:

\begin{corollary}\label{cor:faltings}
  For $b = +1$ ($\delta = -1$), only finitely many $c \in \mathbb{Q}$ give a primitive
  period-$3$ orbit of $f_{c,b}$ over~$\mathbb{Q}$.
\end{corollary}
% SOURCE: exp006 CONCLUSIONS: Faltings applies, finitely many

Similarly, the primitive period-$4$ dynatomic curve at $b = +1$ has degree~$8$
and genus $\leq 21$ (smooth upper bound), so Faltings applies there as well.
% SOURCE: exp006 CONCLUSIONS: period-4 degree 8, genus ≤ 21

\subsection{Resolution of the apparent contradiction}

The universal family~\eqref{eq:family} gives period-$3$ orbits for \emph{every}~$b$
but uses a \emph{different~$c$} for each~$b$: namely
$c(b) = -(29b^2 + 38b + 29)/16$.  In the $(b, c)$-plane, this traces a
parabola --- a rational curve of genus~$0$.

Faltings' theorem constrains the \emph{vertical fibers} $\{b = b_0\}$: for
each fixed~$b_0$, the curve~\eqref{eq:primitive3} (or its analog at $b = b_0$)
has high genus and thus finitely many rational points.
But the universal family moves \emph{horizontally} in the $(b,c)$-plane,
sweeping through these fibers one point at a time.

\begin{center}
\begin{tabular}{@{} lll @{}}
  \toprule
  Direction & Genus & Consequence \\
  \midrule
  Fixed $b$, vary $c$ & $\geq 9$ & Finitely many $c$ (Faltings) \\
  Vary $b$, $c = c(b)$ & $0$ & Infinitely many $(b,c)$ pairs \\
  \bottomrule
\end{tabular}
\end{center}

In algebro-geometric terms, the period-$3$ locus forms a fibration
$\pi\colon V \to \mathbb{A}^1_b$ over the $b$-line, whose fiber $\pi^{-1}(b_0)$
is the dynatomic curve $\Phi_3(x_0, c;\, b_0) = 0$ of genus~$\geq 9$.
By Faltings' theorem, each fiber has finitely many rational points.
The universal family~\eqref{eq:family} is a \emph{rational section}
$\sigma\colon \mathbb{A}^1_b \to V$, $b \mapsto (x_0(b),\, c(b))$,
of this fibration --- a genus-$0$ curve that meets each high-genus fiber
in exactly one rational point.
% SOURCE: src/proof_resultant.py Step 6: universal family verified on period-3 locus

Indeed, the resultant that eliminates $x_1, x_2$ from the period-$3$
system~\eqref{eq:p3a}--\eqref{eq:p3c} at $b = 1$ factors as
$R(x_0, c) = (c + x_0^2) \cdot \Phi_3(x_0, c)$,
where $c + x_0^2 = 0$ is the fixed-point curve and $\Phi_3$ is
equation~\eqref{eq:primitive3}.
% SOURCE: src/proof_resultant.py Steps 3-4: resultant computed, exact division verified

Conjecture~\ref{conj:4} implicitly assumed that the exceptional set would be
determined by the vertical fibers.  The universal family shows that the
``landscape'' of rational periodic orbits has a horizontal structure --- a
rational section of the dynatomic fibration --- that the conjecture did not
anticipate.


% ============================================================================
\section{Discussion and Open Questions}
\label{sec:discussion}
% ============================================================================

\subsection{Canonical heights}

For all nine periodic orbits tested (from the universal family, Conjecture~\ref{conj:3}
sharpness examples, and the period-$7$ orbit), the canonical height
$\hat{h}$ vanishes exactly, as expected.
The minimum non-zero canonical height among non-periodic points tested was
$\hat{h} \geq 0.001$, consistent with Conjecture~5 of~\cite{berger2024}
(the Kawaguchi height lower bound).
% SOURCE: exp005 CONCLUSIONS: ĥ=0 for periodic, min non-zero ĥ ≈ 0.001

The Morton--Silverman uniform boundedness conjecture~\cite{mortonsilverman1994}
(Conjecture~2 of~\cite{berger2024}) is also consistent with our observations: at each
fixed $(c, b)$, the number of rational periodic points is small.
% SOURCE: exp005 CONCLUSIONS: Conj 2 consistent

\subsection{Open questions}

\begin{question}\label{q:period4}
  Does there exist a universal family for period~$4$?  That is, for every
  $b \in \mathbb{Q}$ (or all but finitely many), does there exist
  $c \in \mathbb{Q}$ such that $f_{c,b}$ has a primitive period-$4$ orbit
  over~$\mathbb{Q}$?
\end{question}

\begin{question}\label{q:period5}
  Does any $b \in \mathbb{Q}$ admit a primitive period-$5$ orbit over~$\mathbb{Q}$?
  Our computational search found none for any~$b$ tested.
\end{question}

\begin{question}\label{q:periodset}
  For each $b \in \mathbb{Q}$, what is the full set of periods that occur
  over~$\mathbb{Q}$?  Our data suggests this set depends on the sign of
  $\det Df = -b$, but a complete characterization is open.
\end{question}

\begin{question}\label{q:denomproof}
  Can the denominator preservation result (Proposition~\ref{prop:denom})
  be extended to a complete characterization of which pairs $(q, N)$ admit
  rational orbits of period~$N$ with denominator~$q$?
\end{question}


% ============================================================================
\section*{Reproducibility}
% ============================================================================

All code used to discover and verify the results in this paper is available at:
\begin{center}
  \url{https://github.com/guilhermesilveira/henon-conj4-counterexample}
\end{center}

\begin{itemize}
  \item \texttt{src/exp002j-unified-parametric/proof\_full\_chain.py}: SymPy
  verification of the complete logical chain --- defines the map, computes $f$,
  $f^2$, $f^3$ from the map definition, verifies $f^3(P_0) = P_0$ and
  $f(P_0) \neq P_0$, and spot-checks with exact \texttt{Fraction} arithmetic.

  \item \texttt{src/exp002i-odd-delta-family/verify\_one\_line.py}: standalone
  Python~3 script (no external dependencies) that verifies $20$ period-$3$ orbits
  by exact arithmetic.

  \item \texttt{src/exp002j-unified-parametric/unified\_family\_d.py}: SymPy proof
  that the universal period-$3$ family works identically as a polynomial identity.

  \item \texttt{src/proof\_stability.py}: SymPy proof that the universal orbit is
  always a saddle (Remark~\ref{rem:saddle}). Computes $Df^3$ via the chain rule,
  verifies $\det = -b^3$ and $\operatorname{tr} = -(b-1)(35b^2+62b+35)/8$,
  proves the discriminant quartic $Q(b)$ has no real roots, and performs numeric
  spot-checks.

  \item \texttt{src/proof\_resultant.py}: SymPy proof of the resultant factorization
  $R(x_0,c) = (c+x_0^2)\cdot\Phi_3(x_0,c)$ (Section~\ref{sec:genus}).
  Eliminates $x_1, x_2$ via iterated resultants, verifies exact polynomial division,
  confirms $\Phi_3$ matches equation~\eqref{eq:primitive3}, and checks the
  singularity at $(0,1)$ giving genus~$\geq 9$.  Runtime: $< 1$~second.

  \item \texttt{lean/}: Lean~4 formalization of the main results
  (Theorem~\ref{thm:main}, Corollary~\ref{cor:conj4false}, and
  Proposition~\ref{prop:period7}), using Mathlib. The formalization has
  \textbf{zero} \texttt{sorry} (unproved assertions) and \textbf{zero}
  custom axioms. The three orbit identities are closed by \texttt{field\_simp; ring},
  primitivity by \texttt{linarith}, and the period-$7$ orbit by \texttt{native\_decide}
  over~$\mathbb{Q}$.
\end{itemize}


% ============================================================================
\section*{Data Availability Statement}

The code and computational outputs supporting the results of this study are available at
\url{https://github.com/guilhermesilveira/henon-conj4-counterexample}. Additional
materials are available from the author upon reasonable request.


% ============================================================================
\section*{Disclosure Statement}

The author reports no competing interests.


% ============================================================================
\section*{Funding}

This research received no specific grant from any funding agency in the public,
commercial, or not-for-profit sectors.


% ============================================================================
\section*{Acknowledgments}

Claude Opus 4.6 (Anthropic) was used for assistance with \LaTeX\ drafting,
code generation, literature search and reference discovery, and
English-language improvement. All mathematical claims, proofs, and computations
were independently verified by the author using exact rational arithmetic and
symbolic computation (SymPy).

\bibliographystyle{plain}
\bibliography{references}

\end{document}
